\section{Additional proof details}

\subsection{Object-level digit recovery without semantic division}

It is worth isolating the exact algebra needed for
\cref{thm:uniform-code}, because this is the most important new arithmetic
module.

Let \(D=b^i\) and suppose
\[
 C=R+D(a+bU),\qquad R<D,
\]
with \(D>0\).  The repository theorem \leanname{div\_mod\_unique} is stated for
a successor divisor.  Positivity of \(D\) supplies a predecessor \(D'\) with
\(D=S(D')\).  Instantiate uniqueness with candidate quotient \(a+bU\) and
remainder \(R\).  The required decomposition is exactly the displayed
equality, and the remainder bound is \(R<D\).  Hence
\[
  \operatorname{quo}(C,D)=a+bU,
  \qquad
  \operatorname{rem}(C,D)=R.
\]

For the second division, instantiate uniqueness at divisor \(b\), quotient
\(U\), and remainder \(a\):
\[
  a+bU=bU+a,
  \qquad a<b.
\]
Commutativity of addition puts the equation in the library's expected order.
Therefore
\[
 \operatorname{rem}(\operatorname{quo}(C,D),b)=a.
\]

The proof never reasons about Lean's natural-number quotient.  It reasons only
with the object symbols \leanname{quo} and \leanname{rem} and their checked
uniqueness theorem.

\subsection{The split identity}

For a uniform-base code,
\[
 C_L=\sum_{j<L}a_jb^j,
\]
define
\[
 R_i=\sum_{j<i}a_jb^j,
 \qquad
 U_i=\sum_{k<L-i-1}a_{i+1+k}b^k.
\]
Then
\[
 \begin{split}
 C_L
 &=R_i+a_ib^i+
   \sum_{k<L-i-1}a_{i+1+k}b^{i+1+k}\\
 &=R_i+b^i\left(a_i+bU_i\right).
 \end{split}
\]

In Lean, avoid subtraction-heavy indices if they complicate QF arithmetic.
Prove a recursive suffix decomposition instead:
\[
 C_{i+t+1}=C_i+b^i\bigl(a_i+bU_t\bigr),
\]
where \(U_t\) is built by a forward recursion over the \(t\) remaining
entries.  The statement is QF and its step is a ring identity using
\(b^{n+1}=b^n b\).  This formulation aligns better with the existing sum and
power APIs.

\subsection{Existential discharge with de Bruijn variables}

The schematic producer is
\[
 \Gamma\objectvdash
 \exists c\,[c<B(q)\wedge V(c,q)=1].
\]
The elimination branch lives over one extra free variable and shifted context:
\[
 [c<B(q)\wedge V(c,q)=1],\Gamma^{\uparrow}
 \objectvdash
 \bigl(\Sel(q)<B(q)\wedge V(\Sel(q),q)=1\bigr)^{\uparrow}.
\]

The selected term must ignore the new variable.  Concretely, its occurrence of
\(q\) is shifted, and substitution by the existential variable does not alter
it.  The existing lemmas \leanname{Fm.shift\_subst\_up},
\leanname{Syn.subst\_shift}, and the pattern used in
\leanname{Prf.substWeaken} handle this bookkeeping.  A robust implementation
should first prove a helper lemma:

\begin{lstlisting}[language=Lean]
@[simp] theorem shift_selected (P : Package) (q : PTm n) :
  Syn.shift (selected P q) = selected P (Syn.shift q) := by
  simp [selected, selectT, failureBody]
\end{lstlisting}

Then the branch conclusion simplifies to the shift required by
\leanname{Prf.exE}.

\subsection{Bounded universal reflection}

Let \(a(i,q)\) be Boolean and define
\[
 P(0,q)=1,
 \qquad
 P(n+1,q)=P(n,q)a(n,q).
\]
The QF invariant is
\[
 P(n,q)=1\quad\leftrightarrow\quad
 \text{all previously exposed factors are one}.
\]
The right side should not be placed literally inside the induction matrix.
Instead prove two object lemmas:

\begin{enumerate}
\item \(P(N,q)=1\) and \(i<N\) imply \(a(i,q)=1\), by a zero/product or
factor-extraction theorem;
\item a universally quantified bounded premise that each factor is one implies
\(P(N,q)=1\), using QF induction on the running product.
\end{enumerate}

This is analogous to the existing \leanname{sum\_zero\_elim} and
\leanname{sum\_zero\_of\_bounded} results.  A product version should be added to
the shared core.

\subsection{Master-certificate soundness}

Let the graph checker be
\[
 V_{\mathrm{all}}(C,q)
 =\prod_{i<m}V_i(c_i(C,q),\vec c_{\operatorname{par}(i)}(C,q),q).
\]
Assume it returns one.  Product reflection gives every factor equal to one.
Local checker soundness yields every node contract.  Induct over the fixed node
order and use local extractors to derive the node formulas.  At the terminal
node, obtain the requested theorem.

This gives two formalization styles:

\begin{itemize}
\item prove all local formulas from the global checker by a metalevel induction
on the fixed node list;
\item state a generic list of node descriptors and prove the induction inside
Lean.
\end{itemize}

The first is simpler and fully compatible with fixed-program specialization.
The second is worthwhile only after the arithmetic substrate is stable.

\subsection{Complete-list converse}

Suppose a finite list \(L\) is complete for a reflected predicate \(P(c,q)\),
meaning every valid code equals one of its entries.  Let
\[
 H=1+\max L,
\]
with \(\max\varnothing=0\).  Then every valid code is below \(H\).  The proof
uses existential elimination on list membership, bounded list traversal, and
order transitivity.  Thus a complete list yields a numerical bound.  The
converse from a bound is \cref{thm:height-list}.  This equivalence clarifies
exactly what classical ``finiteness'' lacks: neither side follows without an
explicit presentation of finite evidence.

\subsection{Gap uniqueness}

Suppose labels \(i,j<m\) both pass the approximation checker:
\[
 d(a,v_i)\le e,
 \qquad
 d(a,v_j)\le e.
\]
Symmetry and triangle inequality give
\[
 d(v_i,v_j)\le d(v_i,a)+d(a,v_j)\le2e.
\]
The separation premise says \(i\ne j\to d(v_i,v_j)>2e\), so \(i=j\).  In a
natural/rational coding, strict and non-strict inequalities are reflected by
monus/order terms, and the contradiction uses the repository's order lemmas.
The uniqueness proof can therefore be internalized without decidability of an
arbitrary real equality.
